% =====================================================================
%  09-cli.tex - Command-line reference of the two server binaries
% =====================================================================
\chapter{Referencia de la línea de comandos}\label{cap:cli}

El crate \archivo{crates/bin} produce dos binarios. \texttt{despacho-cli}
reúne el servidor HTTP del despacho, la preparación de la base de datos y las
operaciones criptográficas fuera de línea. \texttt{qadra-ca} es el servicio
del operador de la autoridad certificadora. Este capítulo resume sus subcomandos y
muestra los ejemplos de uso más frecuentes. Los campos de la salida JSON y
las transcripciones de \texttt{despacho-cli} están en el Anexo~B del reporte
técnico (\archivo{latex/chapters/anexo-b-cli.tex}). La definición vigente de
subcomandos y opciones es la del código (\archivo{crates/bin/src/cli.rs},
\archivo{crates/bin/src/database_args.rs}, \archivo{crates/bin/src/serve_args.rs}
y \archivo{crates/bin/src/ca/main.rs}) y se consulta con \texttt{--help}.

\section{Convenciones}\label{sec:cli-convenciones}

\begin{itemize}
  \item \textbf{Ayuda.} \texttt{--help} está disponible en todos los niveles
    de ambos binarios; \texttt{--version}, solo en el comando raíz de
    \texttt{despacho-cli}. Desde el código fuente, la ayuda se consulta con
    \texttt{cargo run --bin despacho-cli -- --help}.
  \item \textbf{Flujos.} Los resultados se escriben en la salida estándar y
    los diagnósticos en el flujo de error. El nivel de diagnóstico se ajusta
    con la variable \texttt{RUST\_LOG}.
  \item \textbf{Salida para programas.} La bandera global \texttt{--json} hace
    que cada operación finita de \texttt{despacho-cli} imprima un objeto JSON
    en lugar de texto. No aplica a \texttt{serve}, que es un proceso
    persistente, ni a \texttt{vault decrypt} sin \texttt{--out}, que escribe
    el contenido descifrado.
  \item \textbf{Entorno.} \texttt{despacho-cli} carga el archivo
    \archivo{.env} del directorio de trabajo, o de uno superior, cuando
    existe; la plantilla es
    \archivo{.env.example}. Las variables se describen en el
    \cref{cap:configuracion}.
\end{itemize}

\subsection{Códigos de salida}\label{sec:cli-salida}

\texttt{despacho-cli} termina con código cero cuando la operación concluye y
con código distinto de cero ante cualquier error. También termina con código
distinto de cero, de forma deliberada, cuando el resultado criptográfico es un
rechazo, de modo que un guion pueda bifurcar sin analizar la salida:

\begin{itemize}
  \item \texttt{verify}: el veredicto es \texttt{not valid}.
  \item \texttt{package export}: la evidencia no verifica antes de exportarse.
  \item \texttt{auth verify-password}: la contraseña no coincide.
  \item \texttt{auth totp verify}: el código es rechazado.
  \item \texttt{audit verify-chain}: la cadena está rota; se nombra el índice
    del primer eslabón roto.
\end{itemize}

No se asignan valores numéricos distintos a cada causa: los guiones deben
comparar únicamente contra cero. \texttt{qadra-ca} devuelve
éxito o fallo; si el servicio no puede iniciar, escribe un mensaje genérico en
el flujo de error, sin detalles de configuración (\cref{cap:problemas}).

\section{Subcomandos de \texttt{despacho-cli}}\label{sec:cli-despacho}

\begin{table}[H]
  \centering
  \caption{Árbol de subcomandos de \texttt{despacho-cli}.}\label{tab:cli-despacho}
  \small
  \begin{tabularx}{\textwidth}{L{2.9cm} L{4.6cm} Y}
    \toprule
    \textbf{Subcomando} & \textbf{Acciones} & \textbf{Propósito} \\
    \midrule
    \texttt{serve} & --
      & Inicia la API HTTP con PostgreSQL, Redis y la TSA local. No ejecuta
        migraciones. \\
    \texttt{database}
      & \archivo{migrate}, \archivo{check}, \archivo{import},
        \archivo{invalidate-restored-password-resets},
        \archivo{invalidate-restored-invitations}
      & Aplica el esquema y los permisos del rol de ejecución, valida la base
        completa, importa un almacenamiento local anterior y cancela
        invitaciones o recuperaciones pendientes tras una restauración. \\
    \texttt{credential-trust} & \texttt{publish}
      & Publica la raíz y la CRL de confianza para declaraciones internas;
        requiere conexión administrativa. \\
    \texttt{pki}
      & \archivo{init-ca}, \archivo{issue}, \archivo{issue-request},
        \archivo{revoke}, \archivo{gen-crl}, \archivo{show}
      & Administra la CA interna mediante los guiones de \archivo{pki/}:
        raíz, emisión (con clave generada o desde una solicitud PKCS\#10),
        revocación, CRL e inspección. \\
    \texttt{crypto} & \texttt{hash}
      & Imprime el resumen SHA-256 de un archivo. \\
    \texttt{vault}
      & \archivo{encrypt}, \archivo{decrypt}, \archivo{rotate-kek}
      & Cifra y descifra paquetes con AES-256-GCM y rota la clave de cifrado
        de claves. \\
    \texttt{sign} & --
      & Firma el resumen de un documento y escribe \texttt{<archivo>.sig}. \\
    \texttt{timestamp} & --
      & Solicita un sello de tiempo RFC~3161; \texttt{--mock} usa la TSA
        local. \\
    \texttt{verify} & --
      & Reporte de integridad, firma, estado del certificado y sello de
        tiempo. \\
    \texttt{package} & \texttt{export}
      & Construye el paquete de evidencia (ZIP) con sus instrucciones de
        verificación. \\
    \texttt{auth}
      & \archivo{calibrate}, \archivo{hash-password},
        \archivo{verify-password}, \texttt{totp enroll}, \texttt{totp verify}
      & Mide el costo de Argon2id en el equipo, genera y verifica hashes de
        contraseña e inscribe y verifica un segundo factor TOTP. \\
    \texttt{audit}
      & \archivo{append}, \archivo{verify-chain}, \archivo{show}
      & Opera la bitácora encadenada de archivo del flujo fuera de línea. La
        API usa la cadena almacenada en PostgreSQL. \\
    \bottomrule
  \end{tabularx}
\end{table}

\section{Subcomandos de \texttt{qadra-ca}}\label{sec:cli-qadra-ca}

\texttt{qadra-ca} es el servicio del portal de la CA, con credenciales,
segundo factor y sesiones propias, separadas de las del despacho. Su línea de
comandos no incluye subcomandos de emisión ni de revocación: en la terminal, esas
operaciones corresponden a \texttt{despacho-cli pki} y a los guiones de
\archivo{pki/}. El funcionamiento del portal se describe en el \cref{cap:ca};
la configuración y operación del servicio, en \archivo{docs/ca/README.md} y
\archivo{docs/ca/operations.md}.

\begin{table}[H]
  \centering
  \caption{Formas de invocación de \texttt{qadra-ca}.}\label{tab:cli-qadra-ca}
  \begin{tabularx}{\textwidth}{L{6.2cm} Y}
    \toprule
    \textbf{Invocación} & \textbf{Propósito} \\
    \midrule
    \texttt{qadra-ca -{}-config <CONFIG>}
      & Ejecuta el servicio del operador con un archivo de configuración
        privado, legible solo por su dueño. \\
    \texttt{qadra-ca enroll-operator} \texttt{-{}-setup-config <RUTA>}
      \texttt{-{}-enrollment-out <RUTA>}
      & Inscribe al operador inicial de la CA a partir de un archivo de
        preparación separado. \\
    \texttt{qadra-ca inspect-public}
      & Inspecciona material público en JSON recibido por la entrada
        estándar, sin modificar el estado de la CA. \\
    \bottomrule
  \end{tabularx}
\end{table}

\section{Ejemplos}\label{sec:cli-ejemplos}

Los comandos de evidencia proceden del Anexo~B del reporte técnico y los de base
de datos, de \archivo{docs/database-operations.md} y
\archivo{docs/database-restore-commands.md}. Se omitió el prefijo de ruta del
binario, se dividieron las líneas largas y las rutas son ilustrativas.

\subsection{Preparar la base de datos e iniciar el servidor}
\label{sec:cli-ej-servidor}

La secuencia completa (\texttt{database migrate} con la conexión
administrativa, \texttt{database check} y \texttt{serve} con el rol de
ejecución) se muestra con sus comandos en las
\cref{sec:inst-migrar,sec:inst-serve}; las opciones de \texttt{serve} y sus
valores por defecto, en la \cref{tab:conf-serve}. \texttt{database migrate}
otorga permisos a un rol de ejecución que debe existir previamente: no crea
la cuenta ni decide su contraseña. \texttt{database check} valida la base
completa con el rol de ejecución:

\begin{verbatim}
despacho-cli database check --json
\end{verbatim}

\subsection{Firmar, sellar y verificar un documento}
\label{sec:cli-ej-evidencia}

La secuencia siguiente supone una CA creada con \texttt{pki init-ca}, un
certificado emitido con \texttt{pki issue --cn "Socia Demo"}, la CRL generada
con \texttt{pki gen-crl} y la TSA local preparada (\cref{cap:ca}).

\begin{verbatim}
despacho-cli sign document.txt \
  --cert pki-ca/certs/socia-demo.crt.pem \
  --key pki-ca/private/socia-demo.key.pem

despacho-cli timestamp document.txt --mock --pki-dir pki

despacho-cli verify document.txt \
  --sig document.txt.sig --tsr document.txt.tsr \
  --cert pki-ca/certs/socia-demo.crt.pem \
  --ca pki-ca/ca.crt.pem --crl pki-ca/crl/crl.pem
\end{verbatim}

\texttt{sign} escribe \archivo{document.txt.sig} y \texttt{timestamp},
\archivo{document.txt.tsr}. \texttt{verify} informa el estado de cuatro
componentes (\texttt{integrity}, \texttt{signature}, \texttt{certificate} y
\texttt{timestamp}) y cierra con la línea \texttt{verdict}, cuyo valor es
\texttt{valid} o \texttt{not valid}.

Un componente sin evidencia aportada se reporta como omitido y no cuenta
contra el veredicto. Con un sello verificado, el certificado se juzga en el
instante del sello.

\subsection{Exportar el paquete de evidencia}\label{sec:cli-ej-paquete}

\begin{verbatim}
despacho-cli package export document.txt \
  --sig document.txt.sig --tsr document.txt.tsr \
  --cert pki-ca/certs/socia-demo.crt.pem \
  --ca pki-ca/ca.crt.pem --crl pki-ca/crl/crl.pem \
  --tsa-chain pki-tsa/tsa-chain.pem --out evidencia.zip
\end{verbatim}

Antes de escribir el ZIP, el subcomando ejecuta la misma verificación que
\texttt{verify} y rehúsa exportar si el veredicto no es válido. La
verificación independiente del paquete se explica en el
\cref{cap:verificacion}.

\importante{El paquete es evidencia técnica de un prototipo con CA y TSA
internas. Un veredicto \texttt{valid} describe la consistencia criptográfica
de los artefactos frente a esas autoridades; no acredita validez jurídica.}
